\begin{lemma}[The hierarchy path number lies in the full-tail range]
Fix a finite loopless multigraph with individual edge identities, a
palette of $k$ colors, and a distinguished edge $e_0$.  Let
$(T,\varphi_n)$ be a positive-rung extended Tashkinov tree constructed
from an arbitrary certified input coloring, let $\mathcal B$ be a
hierarchy base under $\varphi_n$, and let $C$ be a hierarchy
counterexample over that base.  Assume that $\varphi_n$ is proper and
leaves exactly $e_0$ uncolored.

Assume, uniformly for every positive-rung maximum ETT on the same graph,
palette, and distinguished edge, with arbitrary certified input and
generating colorings and with smaller rung number, all six lower-rung
assertions of the technical theorem: elementarity; terminal
interchangeability; at every PE slot, unconditional stability before the
support, stability on the preceding ladder tree when the first-appearance
index is greater than one, and disjointness of stored pairs at distinct PE
slots with distinct supports and equal first-appearance index; the terminal
PE component assertion; the defective-endpoint order assertion; and the
mod-coloring assertion.  Assume also that every
same-rung maximum ETT with a good hierarchy having fewer dividers than
$C$ is elementary.

Assume that $C$ is lexicographically minimal among hierarchy
counterexamples whose ETTs have the same rung number as $T$ and whose
hierarchies have the same divider count, allowing the certified input
coloring, generating coloring, ETT, base, and hierarchy to vary.  Finally
assume
\[
 2n+11\leq |\overline\varphi_n(T_n)|,
\]
where $n$ is the rung number and $T_n$ is the terminal ladder tree.

Write $H=C.\mathit{hierarchy}$, let
\[
 A=\noderef{HierarchyStarredStage}
      (H,\operatorname{Fin.last}(H.\mathit{dividerCount})),
 \qquad \ell=|T|-|A|,
\]
and put $r=\noderef{HierarchyPathNumber}(H)$.  Then
\[
                         1\leq r\leq \ell .
\]
In particular, $r$ is bounded by the full-tail length $\ell$; no
equality between these two parameters is asserted.
\end{lemma}
\begin{proof}
Apply \noderef{HierarchyCounterexampleTailAtLeastTwo} to the ETT, base,
counterexample, properness and unique-uncoloredness assumptions, the
complete lower-rung package, the fewer-divider induction hypothesis, the
fixed-rung equal-divider minimality hypothesis, and the large-branch
estimate stated here.  Its conclusion is
\[
                              2\leq\ell.                       \tag{1}
\]
In particular $|A|\leq |T|$, and the definition of truncated subtraction
now gives
\[
                         |A|+\ell=|T|.                         \tag{2}
\]

Unfold \noderef{HierarchyPathNumber}.  Its candidate set is
\[
 S=\{j\in\{1,\ldots,\ell\}:\text{every exact tail edge after }y_j
       \text{ has endpoint pair prescribed by the displayed path}\}.
\]
The index $j=\ell$ lies in this set.  Indeed, (1) gives
$1\leq\ell\leq\ell$, and the path condition at $\ell$ asks about an edge
position $s$ satisfying $|A|+\ell\leq s<|T|$.  Equation (2) makes these
two inequalities inconsistent.  Thus the condition is vacuous exactly
as for the one-vertex terminal path, and $S$ is nonempty.

The nonempty branch of \noderef{HierarchyPathNumber} defines $r$ as the
least member of the finite set $S$.  Membership of this minimum in $S$
gives $1\leq r\leq\ell$ from the defining interval of $S$.  The fallback
branch is not entered because $\ell\in S$.  This proves both asserted
bounds while keeping $r$ and $\ell$ distinct.
\end{proof}
